\chapter{Results and Analysis}\label{chap:results-analysis}
% Confirm the title with your supervisor
%%%%%%%%%%%%%%%%%%%%%%%%%%%%%
\guideresults %% remove this line when complete
%%%%%%%%%%%%%%%%%%%%%%%%%%%%%

%% Remove this guiding text as you work on the chapter.
Each project is different, and therefore the number and name of chapters may vary. Discuss this with your supervisor. However, unless there is a specific reason to deviate, the 'normal' five chapter structure is encouraged.

The role of this chapter differs depending on your project type and direction. Discuss with your supervisor what the most appropriate approach to take is. In general, the two broad approaches can be broken down into:

\begin{itemize}[nosep]
    \item \textbf {Analysis}-based projects need to detail the process of data collection, processing, and analysis. Where did you collect data? How did you collect it? How did you store and process your data? An analysis is about how results were produced. The interpretation of these results is dealt with in the following Results chapter.
    \item \textbf {Implementation}-driven projects usually aim towards the production of an `artefact' - be this code, a dataset, or a physical object such as a game. This chapter needs to discuss both the design process followed (why did you make it like this?) and the resulting implementation approach (how did you build it?). Any design/implementation trade-offs made should be discussed.
\end{itemize}

